% GENERATED FILE. Edit canonical lessons, then run npm run generate:books.
% canonical-source-hash: fnv1a64:56803ad9451ab10d
% canonical-lessons: SA-C32-tree, SA-C32-leaf, SA-C32-grass, SA-C32-vine, SA-C32-wood

\chapter{The Tree and What Grows}
\label{ch:sa-tree}
\hlchaptermodality{32}

\begin{chapteropening}
\textbf{By the end of this chapter:} \emph{I can name a tree and four things that grow in Sanskrit, and name the two English words that are cousins of grass and of wood.}

Complete SA-C32-wood and demonstrate the chapter promise.
\end{chapteropening}

\section[taruḥ]{\sk{तरुः} (taruḥ) — a tree}

\label{lesson:SA-C32-tree}



\begin{quote}
Before the new word: what did \emph{himam} mean?
\end{quote}

\subsection*{You'll want to know: \sk{तरुः}}
\textbf{\sk{तरुः}} (\emph{taruḥ}) — \textquotedblleft{}a tree\textquotedblright{}.

You have had \textbf{\sk{वनम्}}, a forest, for a long while now. \textbf{\sk{तरुः}} is one tree standing in it.

Sanskrit keeps several names for a tree — \emph{vṛkṣa}, \emph{druma}, and this one — and \textbf{\sk{तरुः}} is the shortest and the easiest to say. It ends in \emph{-uḥ}, like \textbf{\sk{गुरुः}} and \textbf{\sk{वायुः}}, so you already know how it behaves.

\begin{grammarlens}[title={What you've built}]
One tree, picked out of a forest you already know.
\end{grammarlens}

\subsection*{Your turn}
Say these aloud:

\begin{itemize}
\raggedright
  \item \emph{taruḥ}
  \item it once more, slowly
  \item \emph{himam}, then \emph{taruḥ}
\end{itemize}

\begin{tcolorbox}[breakable,colback=teal!4,colframe=teal!35!black,title={Before you move on}]
What does \sk{तरुः} mean? (\textquotedblleft{}a tree\textquotedblright{}.) Say it once more.
\end{tcolorbox}
\section[patram]{\sk{पत्रम्} (patram) — a leaf}

\label{lesson:SA-C32-leaf}



\begin{quote}
Before the new word: what did \emph{taruḥ} mean?
\end{quote}

\subsection*{You'll want to know: \sk{पत्रम्}}
\textbf{\sk{पत्रम्}} (\emph{patram}) — \textquotedblleft{}a leaf\textquotedblright{}.

\textbf{\sk{पत्रम्}} is a leaf. It is also a wing, and also a written page — one word for all three, because the root under it means \textquotedblleft{}to fly\textquotedblright{}. A leaf is the flying thing on a branch; a page is a leaf of another kind.

Greek \emph{pteron}, \textquotedblleft{}wing\textquotedblright{}, is its cousin, behind \emph{helicopter} and \emph{pterodactyl}; so is Latin \emph{penna}, \textquotedblleft{}feather\textquotedblright{}, behind \emph{pen} and \emph{pinnate}; and so is English \emph{feather} itself. Neuter \textbf{-\sk{अम्}}, like \textbf{\sk{वस्त्रम्}} and \textbf{\sk{आसनम्}}.

\begin{grammarlens}[title={What you've built}]
A leaf, a wing and a page, all carried by one word.
\end{grammarlens}

\subsection*{Your turn}
Say these aloud:

\begin{itemize}
\raggedright
  \item \emph{patram}
  \item it once more, slowly
  \item \emph{taruḥ}, then \emph{patram}
\end{itemize}

\begin{tcolorbox}[breakable,colback=teal!4,colframe=teal!35!black,title={Before you move on}]
What does \sk{पत्रम्} mean? (\textquotedblleft{}a leaf\textquotedblright{}.) Say it once more.
\end{tcolorbox}
\section[tṛṇam]{\sk{तृणम्} (tṛṇam) — grass, a blade of grass}

\label{lesson:SA-C32-grass}



\begin{quote}
Before the new word: what did \emph{patram} mean?
\end{quote}

\subsection*{You'll want to know: \sk{तृणम्}}
\textbf{\sk{तृणम्}} (\emph{tṛṇam}) — \textquotedblleft{}grass, a blade of grass\textquotedblright{}.

\textbf{\sk{तृणम्}} covers a single blade and a whole field of it. It is the low green stuff a \textbf{\sk{तरुः}} stands in.

English \emph{thorn} is the same word: Old English \emph{thorn}, Gothic \emph{thaurnus}, from a shape like \textbf{*trnus}, \textquotedblleft{}a stiff thing that grows up out of the ground\textquotedblright{}. Sanskrit kept the mild end of that meaning and Germanic kept the sharp end, which is a common way for one word to become two things.

\begin{grammarlens}[title={What you've built}]
A third growing thing, and the reason grass and a thorn are the same word.
\end{grammarlens}

\subsection*{Your turn}
Say these aloud:

\begin{itemize}
\raggedright
  \item \emph{tṛṇam}
  \item it once more, slowly
  \item \emph{patram}, then \emph{tṛṇam}
\end{itemize}

\begin{tcolorbox}[breakable,colback=teal!4,colframe=teal!35!black,title={Before you move on}]
What does \sk{तृणम्} mean? (\textquotedblleft{}grass, a blade of grass\textquotedblright{}.) Say it once more.
\end{tcolorbox}
\section[latā]{\sk{लता} (latā) — a creeper, a climbing plant}

\label{lesson:SA-C32-vine}



\begin{quote}
Before the new word: what did \emph{tṛṇam} mean?
\end{quote}

\subsection*{You'll want to know: \sk{लता}}
\textbf{\sk{लता}} (\emph{latā}) — \textquotedblleft{}a creeper, a climbing plant\textquotedblright{}.

\textbf{\sk{लता}} is anything that climbs rather than stands: a vine, a creeper, a trailing branch.

Sanskrit poetry pairs it with \textbf{\sk{तरुः}} so often that the two almost arrive together — the creeper and the tree it climbs. Unlike most of this chapter it has no European cousin worth naming; it is homegrown, and it went on to become one of the commonest given names in India. Ends in \textbf{-\sk{आ}}, like \textbf{\sk{तारा}}.

\begin{grammarlens}[title={What you've built}]
A fourth growing thing, and a name half of India answers to.
\end{grammarlens}

\subsection*{Your turn}
Say these aloud:

\begin{itemize}
\raggedright
  \item \emph{latā}
  \item it once more, slowly
  \item \emph{taruḥ}, then \emph{latā}, in the order they grow
\end{itemize}

\begin{tcolorbox}[breakable,colback=teal!4,colframe=teal!35!black,title={Before you move on}]
What does \sk{लता} mean? (\textquotedblleft{}a creeper, a climbing plant\textquotedblright{}.) Say it once more.
\end{tcolorbox}
\section[dāru]{\sk{दारु} (dāru) — wood, timber}

\label{lesson:SA-C32-wood}



\begin{quote}
Before the last word of the chapter: what did \emph{latā} mean, and what did \emph{taruḥ} mean?
\end{quote}

\subsection*{You'll want to know: \sk{दारु}}
\textbf{\sk{दारु}} (\emph{dāru}) — \textquotedblleft{}wood, timber\textquotedblright{}.

\textbf{\sk{दारु}} is the substance a \textbf{\sk{तरुः}} is made of.

It is one of the securest cousins in this book. English \emph{tree} is it, and so are Gothic \emph{triu}, Greek \emph{drus} \textquotedblleft{}oak\textquotedblright{}, Welsh \emph{derwen} and Russian \emph{derevo}; through a Celtic word for one who knows the oak comes \emph{druid}. English \emph{true} and \emph{trust} belong to the same family, on the old idea that to be true is to be as firm as a tree. Ends in short \emph{-u}, like \textbf{\sk{साधु}}.

\begin{grammarlens}[title={What you've built}]
Five growing things, and the tree standing inside the English word \emph{true}.
\end{grammarlens}

\subsection*{Your turn}
Say these aloud:

\begin{itemize}
\raggedright
  \item \emph{dāru}
  \item it once more, slowly
  \item all five in order, then name the English word hiding in \emph{dāru}
\end{itemize}

\begin{tcolorbox}[breakable,colback=teal!4,colframe=teal!35!black,title={Before you move on}]
What does \sk{दारु} mean? (\textquotedblleft{}wood, timber\textquotedblright{}.) Say it once more.
\end{tcolorbox}
